\subsection{The explicit continuum threshold}

The high-rate shape theorem determines the threshold, including its
correct algebraic branch. Write $\rho_*=8-3\sqrt6$.
\begin{theorem}[Sharp continuum threshold above half rate]
For $1/2<\rho\le\rho_*$, the threshold is $a_{\rm DKT}(\rho)$.
For $\rho_*<\rho<1$, put
\[
 p=\sqrt{2\rho-1},\qquad \delta=1-p,\qquad
 \gamma=\sqrt\delta\,\frac{(1+p)(2-p)}{1+p^2}.
\]
Let $y$ be the largest real root of $y^3-3y+\gamma=0$. Then the threshold is
\begin{equation}\label{eq:high-threshold}
 B_* =\frac{p\sqrt\delta}{(1+p^2)y-\delta^{3/2}},
 \qquad a_* =\frac{1+p}{2}+\frac{\delta^2}{2}B_*.
\end{equation}
It satisfies $a_*<a_{\rm DKT}(\rho)$. At the threshold a nonempty
support has zero surplus; below it every positive-area support has
negative surplus, and above it positive surplus is possible.
These are continuum assertions, not finite-multiplicity converses
at high rates.
\end{theorem}
\begin{proof}
First set $a=a_{\rm DKT}(\rho)$ and
$B_0=3(1-a)/(2(2-\rho))$. The saturated cap polynomial and its derivative
factor as
\[
 P(B)=-\frac{2-\rho}{6}B(B-B_0)^2,
 \qquad f(s)=-\frac{2-\rho}{2}(s-B_0/3)(s-B_0).
\]
The substitution calculation in the preceding corollary also proves
$c\ge B_0$ when $1/2<\rho\le\rho_*$. In that case $P(B)\le0$ for
$B\le c$. On the remaining interval $[c,d]$, the optimized column
surplus $M$ of \eqref{eq:tail} is convex, with
$M(c)=f(c)\le0$ and $M(d)<0$. It is therefore nonpositive throughout
that interval. No width gives positive surplus, while $B=B_0$ gives
zero. Increasing agreement strictly improves the benefit of this
nonempty support. A zero-surplus positive-area support at any smaller
agreement would similarly give positive surplus before this threshold,
a contradiction. This proves the first case.

Now suppose $\rho>\rho_*$. At $a_{\rm low}=(1+p)/2$ one has
$d=p/(1+p)$. For $0<s\le d$,
\[
 d-s<p(1-d-s),
\]
so $M(s)<0$; the full-column surplus satisfies $f(s)\le M(s)$.
Thus every nonempty optimized prefix has negative surplus there.
There is positive surplus below $a_{\rm DKT}$ by the preceding
corollary. The derivative of total surplus at $B=0$ is
$a^2/(2\rho)-1/2$, uniformly negative on
$[a_{\rm low},a_{\rm DKT}]$. Widths approaching zero therefore cannot
supply the first positive surplus. Continuity and compactness give a
first threshold strictly between these two agreements, with a
nonempty zero-surplus maximizing width $B$.

At this threshold, $B>c$: below $a_{\rm DKT}$ the saturated polynomial
$P(B)$ is strictly negative for every $B>0$. Also $B<d$, since
$M(d)<0$ would make a slightly smaller width have positive surplus if
$B=d$ had zero surplus. Thus stationarity gives $M(B)=0$.
Both $d-B$ and $1-d-B$ are positive, so this equation has the definite
sign choice
\[
 d-B=p(1-d-B),\qquad
 a=\frac{1+p}{2}+\frac{(1-p)^2}{2}B.
\]
Put $x=c/B$. Then
\[
 B=\frac{p\delta}{(1+p^2)x-\delta^2},\qquad
 Q(x)=x^3-3\delta x+
        \delta^2\frac{(1+p)(2-p)}{1+p^2}.
\]
Substitution into the total surplus of \eqref{eq:high-total} gives exactly
\[
 -\frac{p\delta^3(1+p^2)}
 {12\bigl((1+p^2)x-\delta^2\bigr)^3}\,Q(x).
\]
Hence $Q(x)=0$.

For the branch choice, $B<d<1/2$ implies
$x>x_{1/2}=(1-p^2)/(1+p^2)$, and
\[
 Q(0)>0,\qquad
 Q(x_{1/2})=-\frac{p^2(1-p)^2(1+p)^2(p^2+3)}{(1+p^2)^3}<0.
\]
The cubic has two positive roots, exactly one greater than $x_{1/2}$;
its third root is negative. Thus the threshold selects the largest
root. Setting $x=\sqrt\delta\,y$ proves
\eqref{eq:high-threshold}. Existence and admissibility of this root
follow from the interior optimizer just established, rather than from
an unsupported choice of a numerical cubic branch.

Finally, the threshold support gains strictly when agreement increases.
A positive-area zero-surplus support below threshold would also gain
strictly at a slightly higher agreement still below threshold, which
is impossible. This proves the stated interpretation on both sides.
\end{proof}

The formula and proof are recovered from the September 13 draft, with
an expanded branch and existence audit here. Exact symbolic checks
verify the three algebraic identities used above. Rational isolation
at six rational values of $p$ produces admissible root intervals and
agreement intervals of width less than $10^{-28}$, with opposite
surplus signs at their endpoints. These checks supplement the proof;
they are not a substitute for the quantified threshold statement.
